\section{Method}\label{sec:method}
For each position $p$ and each angular bin (binned in sky tangents in world coordinates) we define the opacity
\begin{equation}
\begin{aligned}
\lambda_p&=-\ln\!\left[\frac{n_p}{(t_p/t_{\mathrm{sky}})\,n_{\mathrm{sky}}}\right],\\
\sigma_{\lambda,p}^2&\simeq\frac{1}{n_p}+\frac{1}{n_{\mathrm{sky}}},
\end{aligned}
\label{eq:opacity}
\end{equation}
where $n$ are track counts and $t$ live times. The open-sky ratio normalizes the detector response, assuming it remained stable between runs. The live-time ratio fixes the zero of $\lambda$, so no free offsets are fitted in the nominal voxel inversion. Nominal fit weights use binwise standard deviations estimated from \InputsOpacityWeightReplicas{} Poisson replicas; Eq.~\ref{eq:opacity} gives the analytic approximation. Bins with too few sky tracks or at steep angles are masked, bins whose replicas show no spread get zero weight, and opacities are capped by a transmission floor. Figure~\ref{fig:opacity} shows the opacity maps.

\begin{figure*}[t]
\centering
\includegraphics[width=0.88\textwidth]{opacity.pdf}
\caption{Opacity seen from the two room positions, and open-sky counts. The beam shadows appear as repeated stripes on a smooth background from the slab.}\label{fig:opacity}
\end{figure*}

For the three-dimensional image we write $\boldsymbol\lambda=A\boldsymbol q$, where $q_j$ is the opacity per unit length in voxel $j$ and $A_{ij}$ is the path length of bin $i$ through voxel $j$, averaged over the detector aperture. A track with local detector tangents $(t_x,t_y)$ can only pass through all layers within an area of size $(W-|t_x|D)\times(W-|t_y|D)$, with $W$ the active width and $D$ the layer gap; we sample this area with sub-rays and rotate it with the detector azimuth. The system is solved with non-negative iterative updates \citep{gilbert1972} and total-variation smoothing in the horizontal plane only \citep{rudin1992}. We do not smooth vertically, because the height information comes from parallax and smoothing would blur it. Voxels with low weighted coverage are damped to suppress edge artefacts. The smoothing strength was chosen on the full data, and the image is the iterate with the lowest weighted residual within a fixed budget of \ReconstructionIterations{} updates. The voxel size is $\ReconstructionSpacing\,\mathrm{m}$; this is the grid, not the resolution, which is much coarser in $z$.

The pose of pos1 (floor offset, separation and azimuth) is fitted by requiring that a reconstruction from one position predicts the data of the other. Both positions are assumed to share one floor level. This objective reconstructs each view alone on a coarser lattice, without smoothing and with analytic weights, and allows a free offset on the predicted view. To avoid locking onto the beam pitch, it is averaged over two such lattices shifted horizontally by half a cell. It gives a floor offset $x=\PoseX\pm\PoseXSigma\,\mathrm{m}$, $y=\PoseY\pm\PoseYSigma\,\mathrm{m}$, a separation of $\PoseBaseline\,\mathrm{m}$ and an azimuth of $\PoseAz\pm\PoseAzSigma^\circ$, with count-replica spreads. The azimuth is not resolved: small rotations often leave the binned data unchanged. The absolute scale remains conditional on the rough floor-offset prior $\InputsFloorOffsetPrior\,\mathrm{m}$, search bounds and reconstruction objective. The finite detector footprint and voxel lattice also introduce length scales; the fitted pose is not an independent survey. On-site ceiling measurements were not used to choose it.
